\documentclass[11pt,a4paper]{article}
\usepackage[spanish]{babel}
\usepackage[utf8]{inputenc}
\usepackage[T1]{fontenc}
\usepackage{lmodern}
\usepackage{inconsolata} % mejor fuente monoespaciada para código
\usepackage{geometry}
\geometry{margin=2cm}
\usepackage{microtype}
\usepackage{hyperref}
\hypersetup{colorlinks=true, linkcolor=black, urlcolor=black, citecolor=black}
\usepackage[dvipsnames]{xcolor}
\usepackage{enumitem}
\usepackage{graphicx}
\usepackage{tikz}
\usepackage{tcolorbox}
\tcbuselibrary{skins,breakable, listings, theorems}
\usepackage{titlesec}
\usepackage{listings}
\usepackage{listingsutf8}
\usepackage{booktabs}
\usepackage{caption}
% Estilo y navegación del índice/ToC
\usepackage{bookmark}
\usepackage{tocloft}
\usepackage{fancyhdr}
\usepackage{parskip}

% Colores de enlaces más acordes al tema
\hypersetup{colorlinks=true, linkcolor=black, urlcolor=black, citecolor=black}

% Título del índice en español
\addto\captionsspanish{\renewcommand{\contentsname}{Índice}}

% Profundidad del índice y numeración
\setcounter{tocdepth}{2}
\setcounter{secnumdepth}{3}

% Estilo del índice (ToC)
\renewcommand{\cfttoctitlefont}{\LARGE\bfseries}
\renewcommand{\cftaftertoctitle}{\par\vspace{6pt}\hrule\vspace{10pt}}
\renewcommand{\cftsecleader}{\cftdotfill{\cftdotsep}}
\renewcommand{\cftsecfont}{\bfseries}
\renewcommand{\cftsecpagefont}{\bfseries}
\setlength{\cftbeforesecskip}{4pt}
\setlength{\cftbeforesubsecskip}{2pt}

% Estilo de títulos
\titleformat{\section}{\Large\bfseries\color{accent}}{\thesection}{0.6em}{}[\vspace{2pt}\titlerule]
\titlespacing*{\section}{0pt}{2.0ex plus 1ex minus .2ex}{1ex}
\titleformat{\subsection}{\large\bfseries\color{accent2}}{\thesubsection}{0.6em}{}
\titlespacing*{\subsection}{0pt}{1.6ex plus .8ex minus .2ex}{0.8ex}
\titleformat{\subsubsection}{\normalsize\bfseries}{\thesubsubsection}{0.6em}{}

% Encabezados y pies de página
\pagestyle{fancy}
\fancyhf{}
\fancyhead[L]{Angular a Nivel Senior}
\fancyhead[R]{\nouppercase{\leftmark}}
\fancyfoot[C]{\thepage}
\renewcommand{\headrulewidth}{0.4pt}
\renewcommand{\footrulewidth}{0pt}
\setlength{\headheight}{14pt}

% ------------------------------------------------------
% Paquetes básicos
% ------------------------------------------------------
\usepackage{titlesec}
\usepackage{listings}
\usepackage{listingsutf8}
\usepackage{booktabs}
\usepackage{caption}
\captionsetup[lstlisting]{labelfont=bf, singlelinecheck=false, margin=0pt}

% Estilos visuales
\definecolor{codebg}{HTML}{F5F7FA}
\definecolor{codefg}{HTML}{111827}
\definecolor{codekw}{HTML}{7C3AED}
\definecolor{codestr}{HTML}{8B4513}
\definecolor{codecm}{HTML}{6B7280}
\definecolor{codeid}{HTML}{111827}
\definecolor{accent}{HTML}{111827}
\definecolor{accent2}{HTML}{374151}

\lstdefinelanguage{TypeScript}{
  sensitive=true,
  morekeywords={abstract,as,asserts,any,async,await,boolean,break,case,catch,class,const,continue,debugger,declare,default,delete,do,else,enum,export,extends,false,finally,for,from,function,get,if,implements,import,in,infer,instanceof,interface,is,keyof,let,module,namespace,never,new,null,number,object,of,package,private,protected,public,readonly,require,global,return,set,string,symbol,throw,true,try,type,typeof,undefined,unique,unknown,var,void,while,yield},
  morecomment=[l]{//},
  morecomment=[s]{/*}{*/},
  morestring=[b]',
  morestring=[b]",
}

\lstdefinelanguage{HTML}{
  sensitive=true,
  morecomment=[s]{<!--}{-->},
  morestring=[b]',
  morestring=[b]",
  alsoletter={<>=-},
  morekeywords={@if,@for,@switch,@case,@default,div,span,section,article,header,footer,main,nav,button,input,form,ng-container,router-outlet},
}

\lstdefinestyle{tsstyle}{
  language=TypeScript,
  basicstyle=\ttfamily\small\color{codefg},
  backgroundcolor=\color{codebg},
  keywordstyle=\bfseries\color{codekw},
  stringstyle=\color{codestr},
  commentstyle=\itshape\color{codecm},
  identifierstyle=\color{codefg},
  numberstyle=\tiny\color{codecm},
  numbers=left,
  stepnumber=1,
  numbersep=8pt,
  showstringspaces=false,
  tabsize=2,
  breaklines=true,
  frame=single,
  rulecolor=\color{black!15},
  frameround=tttt,
  framesep=6pt,
  xleftmargin=2pt,
  xrightmargin=2pt
}

\lstdefinestyle{htmlstyle}{
  language=HTML,
  basicstyle=\ttfamily\small\color{codefg},
  backgroundcolor=\color{codebg},
  keywordstyle=\color{codekw}\bfseries,
  stringstyle=\color{codestr},
  commentstyle=\itshape\color{codecm},
  identifierstyle=\color{codefg},
  numberstyle=\tiny\color{codecm},
  numbers=left,
  stepnumber=1,
  numbersep=8pt,
  showstringspaces=false,
  tabsize=2,
  breaklines=true,
  frame=single,
  rulecolor=\color{black!15},
  frameround=tttt,
  framesep=6pt,
  xleftmargin=2pt,
  xrightmargin=2pt
}

% Habilitar UTF-8 en listings y mapear acentos/ellipsis para pdflatex
\lstset{
  inputencoding=utf8,
  literate={á}{{\'a}}1 {é}{{\'e}}1 {í}{{\'\i}}1 {ó}{{\'o}}1 {ú}{{\'u}}1
           {Á}{{\'A}}1 {É}{{\'E}}1 {Í}{{\'I}}1 {Ó}{{\'O}}1 {Ú}{{\'U}}1
           {ñ}{{\~n}}1 {Ñ}{{\~N}}1
           {ü}{{\"u}}1 {Ü}{{\"U}}1
           {…}{{\ldots}}1
}

\newtcolorbox{tipbox}{enhanced, breakable, colback=accent!5, colframe=accent!60!black, coltitle=black, title=Consejo Senior, attach boxed title to top left={yshift=-2mm, xshift=2mm}, boxed title style={colback=accent!20, sharp corners}}
\newtcolorbox{warnbox}{enhanced, breakable, colback=orange!5, colframe=orange!70!black, coltitle=black, title=Evita esto, attach boxed title to top left={yshift=-2mm, xshift=2mm}, boxed title style={colback=orange!20, sharp corners}}
\newtcolorbox{refbox}{enhanced, breakable, colback=green!5, colframe=green!60!black, coltitle=black, title=Patrón, attach boxed title to top left={yshift=-2mm, xshift=2mm}, boxed title style={colback=green!20, sharp corners}}

\title{Angular a Nivel Senior: Guía Práctica y de Referencia}
\author{Apuntes de trabajo}
\date{\today}

\begin{document}
\maketitle
\tableofcontents
\newpage

\section{Introducción}
Estos apuntes condensan prácticas, patrones y técnicas de Angular modernas (v16+), con foco en componentes standalone, reactividad con Signals, control de flujo de plantillas (@if, @for), vistas diferidas (@defer), DI con \texttt{inject()}, routing funcional y optimización de rendimiento (OnPush, zoneless, SSR + hydration). Se asume solidez en TypeScript.

\begin{tipbox}
Apunta a: simplicidad, composición, tipado estricto y límites claros entre UI, estado y efectos secundarios.
\end{tipbox}

% --- NUEVO: sección introductoria paso a paso para nivel más bajo ---
\section{Fundamentos prácticos: de cero a tu primer componente}
Esta sección explica Angular desde cero con ejemplos breves y lenguaje simple. Si ya tienes experiencia, salta a las siguientes secciones.

\subsection{Antes de Angular: cómo funciona la web}
\begin{itemize}[leftmargin=*]
  \item HTML estructura el contenido (\texttt{<h1>}, \texttt{<p>}, \texttt{<button>}).
  \item CSS da estilos (colores, tamaños, layouts).
  \item JavaScript añade comportamiento (responder a clicks, pedir datos al servidor).
\end{itemize}
Angular es un framework que organiza todo esto para construir aplicaciones grandes de forma mantenible.

\subsection{Qué es un componente}
Un componente es una pieza de UI con: (1) plantilla (HTML), (2) lógica (TypeScript) y (3) estilos. Reutilizas componentes como \emph{bloques Lego}.

\begin{lstlisting}[style=tsstyle, caption={Componente mínimo standalone}]
import { Component } from '@angular/core';

@Component({
  selector: 'app-hello',           // etiqueta HTML que usarás
  standalone: true,                // no necesita NgModule
  template: `<h1>Hola Angular</h1>`
})
export class HelloComponent {}
\end{lstlisting}

\subsection{Data binding: conecta datos con la vista}
\begin{itemize}[leftmargin=*]
  \item Interpolación: muestra valores con \texttt{{\{}\{ valor \}\}}.
  \item Property binding: \texttt{[prop]} pasa datos a atributos/propiedades.
  \item Event binding: \texttt{(evento)} escucha acciones del usuario.
\end{itemize}

\begin{lstlisting}[style=tsstyle, caption={Estado local sencillo con Signals}]
import { Component, signal } from '@angular/core';

@Component({
  selector: 'app-counter',
  standalone: true,
  template: `
    <h2>Contador: {{ count() }}</h2>
    <button (click)="inc()">+1</button>
    <button (click)="reset()" [disabled]="count() === 0">Reset</button>
  `
})
export class CounterComponent {
  count = signal(0);              // estado local
  inc() { this.count.update(c => c + 1); }
  reset() { this.count.set(0); }
}
\end{lstlisting}

\subsection{Control de flujo en la plantilla}
Muestra/oculta y repite elementos sin lógica compleja en TS.
\begin{lstlisting}[style=htmlstyle, caption={@if y @for con trackBy básico}]
<h3>Tareas</h3>
@if (todos().length === 0) {
  <p>No hay tareas</p>
} @else {
  <ul>
    @for (t of todos(); track t.id) {
      <li>{{ t.title }}</li>
    }
  </ul>
}
\end{lstlisting}

\subsection{Servicios e inyección de dependencias (DI)}
Un servicio encapsula lógica reutilizable (por ejemplo, hablar con una API). Angular te lo “inyecta” donde lo necesites.
\begin{lstlisting}[style=tsstyle, caption={Servicio sencillo e inyección con inject()}]
import { Injectable, inject, Component } from '@angular/core';

@Injectable({ providedIn: 'root' })
export class GreeterService {
  greet(name: string) { return `Hola, ${name}!`; }
}

@Component({ selector: 'app-greeter', standalone: true, template: `<p>{{ msg }}</p>` })
export class GreeterComponent {
  private greeter = inject(GreeterService);
  msg = this.greeter.greet('Mundo');
}
\end{lstlisting}

\subsection{HTTP básico}
\begin{lstlisting}[style=tsstyle, caption={Pedir datos con HttpClient (GET)}]
import { HttpClient } from '@angular/common/http';
import { inject, Component, signal } from '@angular/core';

@Component({ selector: 'app-users', standalone: true, template: `
  <button (click)="load()">Cargar</button>
  <ul>@for (u of users(); track u.id) { <li>{{ u.name }}</li> }</ul>
` })
export class UsersComponent {
  private http = inject(HttpClient);
  users = signal<ReadonlyArray<{id:number; name:string}>>([]);
  load() {
    this.http.get<{id:number; name:string}[]>('https://jsonplaceholder.typicode.com/users')
      .subscribe(xs => this.users.set(xs));
  }
}
\end{lstlisting}

\subsection{Routing básico}
Define rutas (URL \textrightarrow{} componente) y navega entre pantallas.
\begin{lstlisting}[style=tsstyle, caption={Rutas mínimas con loadComponent}]
import { Routes } from '@angular/router';

export const routes: Routes = [
  { path: '', loadComponent: () => import('./home.component').then(m => m.HomeComponent) },
  { path: 'about', loadComponent: () => import('./about.component').then(m => m.AboutComponent) },
  { path: '**', redirectTo: '' }
];
\end{lstlisting}

\begin{tipbox}
Sigue este orden para aprender: componentes \textrightarrow{} data binding \textrightarrow{} servicios/DI \textrightarrow{} HTTP \textrightarrow{} routing. Cuando lo domines, vuelve a la sección de \textbf{Signals}, \textbf{@defer} y \textbf{OnPush} para optimizar.
\end{tipbox}

% --- FIN nuevo bloque de fundamentos ---

% --- NUEVO: Glosario ultra básico ---
\section{Glosario ultra básico (palabras difíciles, explicado simple)}
\begin{itemize}[leftmargin=*]
  \item \textbf{Arranque (bootstrap)}: “encender” la app. En Angular se hace en \texttt{main.ts} con \texttt{bootstrapApplication(AppComponent, { providers: [...] })}. No tiene relación con \textbf{Bootstrap} (framework de CSS).
  \item \textbf{Componente}: pieza de la pantalla con su HTML (plantilla) y lógica (TypeScript). Piensa en bloques de Lego.
  \item \textbf{Plantilla (template)}: el HTML que renderiza el componente. Puede contener expresiones como \texttt{{\{\{ titulo \}\}}} y control de flujo \texttt{@if/@for}.
  \item \textbf{Estado}: datos que cambian con el tiempo (contador, lista). Con \textbf{Signals} se guarda con \texttt{signal()} y se deriva con \texttt{computed()}.
  \item \textbf{Servicio (service)}: clase para lógica reutilizable (hablar con API, cache). No dibuja UI.
  \item \textbf{Inyección de dependencias (DI)}: Angular te “entrega” servicios donde los pides con \texttt{inject(MiServicio)}.
  \item \textbf{Provider}: declaración que dice “para este token, usa esta implementación/valor”. Sirve para registrar servicios/configuraciones.
  \item \textbf{Ruta (route)/Router}: relación URL \textrightarrow{} componente. Permite navegar entre pantallas sin recargar la página.
  \item \textbf{HttpClient}: utilidad para hacer \texttt{GET/POST/...} a APIs.
  \item \textbf{Interceptor}: pieza que envuelve cada petición HTTP (añadir token, reintentos, cache).
  \item \textbf{Guard}: función que decide si se puede entrar a una ruta (por ejemplo, si hay sesión).
  \item \textbf{Resolver}: trae datos antes de cargar una ruta.
  \item \textbf{Standalone}: componentes/servicios que no necesitan \texttt{NgModule}. Es el estilo moderno por defecto.
  \item \textbf{@if / @for}: nuevo control de flujo en plantillas. Más claro que \texttt{*ngIf/*ngFor} clásicos.
  \item \textbf{Lazy loading}: cargar características (código) sólo cuando se navega a ellas. Mejora el arranque.
  \item \textbf{SSR (Server-Side Rendering)}: renderizar HTML en el servidor para entregar contenido rápido y amigable a buscadores.
  \item \textbf{Hydration}: al llegar al navegador, Angular conecta eventos sobre el HTML ya renderizado por SSR.
  \item \textbf{Detección de cambios (OnPush)}: estrategia para renderizar menos y más rápido. No mutar objetos “en sitio”.
  \item \textbf{Zone.js / Zoneless}: mecanismo clásico para saber cuándo actualizar la UI (Zone.js). En “zoneless”, tú controlas con Signals y eventos.
  \item \textbf{Observable (RxJS)}: flujo de valores en el tiempo (como un radio). Útil para peticiones, eventos, websockets.
  \item \textbf{NgRx/Store}: patrón para estado global complejo y trazable (acciones, reducers, selectors).
  \item \textbf{Build/Bundle}: proceso de empaquetar todo el código para producción (minificado, dividido en chunks).
\end{itemize}

\begin{warnbox}
\textbf{Bootstrap (arranque)} \neq \textbf{Bootstrap (CSS)}. Aquí “bootstrap” significa “iniciar la app”.
\end{warnbox}

% --- NUEVO: Tutorial guiado ---
\section{Tutorial guiado: de “Hola” a una mini app de tareas}
Meta: pasar de un \emph{hola} a una lista de tareas simple usando los conceptos clave, sin trucos.

\subsection{Paso 1: Componente raíz simple}
Crea un componente standalone con un título y un párrafo. Revisa el ejemplo de \texttt{HelloComponent} en Fundamentos.

\subsection{Paso 2: Estado local con Signals}
Agrega un contador con \texttt{signal(0)}, un botón que incrementa y otro que resetea. Ya tienes un dato que cambia y se refleja en la vista.

\subsection{Paso 3: Lista de tareas (estado compuesto)}
Define un tipo \texttt{Todo} y guarda una lista en \texttt{signal<ReadonlyArray<Todo>>([])}. Añade funciones \texttt{add(title)}, \texttt{toggle(id)} y \texttt{remove(id)} usando copias (no mutaciones).

\begin{lstlisting}[style=tsstyle, caption={Mini store local de Tareas con Signals}]
export type Todo = Readonly<{ id: string; title: string; done: boolean }>; 

export class TodosComponent {
  readonly todos = signal<ReadonlyArray<Todo>>([]);
  readonly remaining = computed(() => this.todos().filter(t => !t.done).length);

  add(title: string) {
    const todo: Todo = { id: crypto.randomUUID(), title, done: false };
    this.todos.update(list => [todo, ...list]);
  }
  toggle(id: string) {
    this.todos.update(list => list.map(t => t.id === id ? { ...t, done: !t.done } : t));
  }
  remove(id: string) {
    this.todos.update(list => list.filter(t => t.id !== id));
  }
}
\end{lstlisting}

\begin{lstlisting}[style=htmlstyle, caption={Plantilla de tareas con @for y @if}]
<h2>Tareas (pendientes: {{ remaining() }})</h2>
<input #i placeholder="Nueva tarea"/>
<button (click)="add(i.value); i.value=''">Añadir</button>

@if (todos().length === 0) { <p>No hay tareas</p> } @else {
  <ul>
    @for (t of todos(); track t.id) {
      <li>
        <label>
          <input type="checkbox" [checked]="t.done" (change)="toggle(t.id)"/>
          <span [class.done]="t.done">{{ t.title }}</span>
        </label>
        <button (click)="remove(t.id)">×</button>
      </li>
    }
  </ul>
}
\end{lstlisting}

\subsection{Paso 4: Separar lógica en un servicio}
Mueve la gestión de tareas a un \texttt{TodosService} y \texttt{inject()} ese servicio en tu componente. Beneficio: testeable y reutilizable.

\subsection{Paso 5: HTTP cuando haga falta}
Sustituye los datos locales por una llamada con \texttt{HttpClient}. Usa \texttt{switchMap} para búsquedas y maneja errores con \texttt{catchError}. No recargues la vista manualmente: las señales u observables lo harán.

\subsection{Paso 6: Routing para varias pantallas}
Crea rutas simples: inicio (lista) y detalle (tarea). Usa \texttt{loadComponent} y, si necesitas datos previos, un \texttt{ResolveFn}.

\begin{tipbox}
Construye siempre en pasos pequeños: primero UI, luego estado local, después servicios, y finalmente HTTP y routing. Así entenderás cada pieza sin abrumarte.
\end{tipbox}

% --- FIN nuevo tutorial guiado ---

\section{Ecosistema Angular actual}
\paragraph{Qué significa y por qué importa}
Angular moderno simplifica la creación de apps grandes: los componentes \emph{standalone} eliminan la fricción de módulos, el nuevo control de flujo (@if/@for) hace las plantillas más legibles, y \textbf{Signals} ofrecen un modelo de estado simple y predecible. \textbf{@defer} y el \emph{lazy loading} reducen el tiempo de carga inicial, mientras que SSR + \emph{hydration} mejoran SEO y rendimiento percibido.

\paragraph{En la práctica}
Piensa en una app como páginas (rutas) compuestas por componentes. Cada componente tiene estado local (Signals), puede pedir datos (HttpClient) y reacciona a eventos de usuario. Divide tu app por \textbf{features} (dominios) y carga lo que hace falta cuando hace falta.

\begin{itemize}[leftmargin=*]
  \item Standalone por defecto: adiós a \texttt{NgModule} en apps nuevas; se configura con providers y \texttt{bootstrapApplication}.
  \item Nuevo control de flujo: \texttt{@if}, \texttt{@for}, \texttt{@switch}.
  \item Reactividad con \textbf{Signals}: \texttt{signal}, \texttt{computed}, \texttt{effect}.
  \item \textbf{@defer}: fragmentos de plantilla diferidos por condición, interacción u ociosidad.
  \item Builder esbuild y servidor rápido; SSR + \emph{hydration} estable.
  \item Formularios tipados; \texttt{takeUntilDestroyed} para limpieza de suscripciones.
\end{itemize}

\begin{warnbox}
Evita mezclar demasiada lógica en plantillas. Mueve cálculos a \texttt{computed()} o a servicios, y reutiliza componentes para mantener la UI simple.
\end{warnbox}

\section{TypeScript para Angular (rápido y efectivo)}
\paragraph{Qué aprenderás aquí}
Modelar datos y contratos claros para que el compilador te ayude: tipos de dominio, DTOs, utilitarios y técnicas de \emph{narrowing} en eventos y HTTP.

\paragraph{Por qué \texttt{type} sobre \texttt{interface}}
\texttt{type} compone mejor (uniones/intersecciones). Usa \texttt{interface} cuando necesites \emph{declaration merging} o describir clases públicas.

\paragraph{Narrowing efectivo}
Evita \texttt{any}. Usa \texttt{typeof}/\texttt{in}/\texttt{instanceof} para guiar al compilador cuando tratas con \texttt{unknown} o eventos del DOM.

\begin{lstlisting}[style=tsstyle, caption={Narrowing en handlers y datos HTTP}]
function onInput(e: Event) {
  const el = e.target as HTMLInputElement; // narrowing seguro
  const value = el.value.trim();
}

http.get<unknown>('/api/user').subscribe(x => {
  if (isUser(x)) { setUser(x); } else { showError(); }
});
\end{lstlisting}

\begin{itemize}[leftmargin=*]
  \item Preferir \texttt{type} sobre \texttt{interface} salvo necesidad de declaración de fusión.
  \item Domina utilitarios: \texttt{Partial}, \texttt{Required}, \texttt{Pick}, \texttt{Omit}, \texttt{Record}, \texttt{ReturnType}.
  \item \textbf{Narrowing} con \texttt{in}, \texttt{typeof}, \texttt{instanceof}. Nunca uses \texttt{any}.
\end{itemize}

\begin{tipbox}
Crea tipos para IDs, DTOs y estado de UI. Te ahorrará bugs y te dará autocompletado útil.
\end{tipbox}

\section{Arquitectura y organización}
\paragraph{Cómo organizar un proyecto}
Organiza por \textbf{features} (dominios), no por tipos de archivo. Cada feature contiene sus componentes, servicios y rutas. \texttt{shared/} para UI reutilizable y \texttt{core/} para servicios globales (auth, http). Con \textbf{standalone}, la mayoría de piezas no requieren \texttt{NgModule}.

\subsection{Standalone y bootstrap}
\paragraph{Arranque en 3 pasos}
1) \texttt{bootstrapApplication(AppComponent)}. 2) Registra \texttt{providers} globales (router, HttpClient, Material, etc.). 3) Añade interceptores/estrategias como funciones puras.

\begin{lstlisting}[style=tsstyle, caption={Arranque con standalone y providers raíz}]
// main.ts
import { bootstrapApplication } from '@angular/platform-browser';
import { provideHttpClient, withInterceptors } from '@angular/common/http';
import { provideRouter } from '@angular/router';
import { AppComponent } from './app/app.component';
import { routes } from './app/app.routes';

bootstrapApplication(AppComponent, {
  providers: [
    provideRouter(routes),
    provideHttpClient(
      withInterceptors([authInterceptor, loggingInterceptor])
    ),
  ]
});
\end{lstlisting}

\subsection{Inyección de dependencias avanzada}
\paragraph{Qué problema resuelve}
Compartir lógica (auth, cache, configuración) sin acoplar componentes a implementaciones concretas. Con \texttt{inject()} puedes obtener servicios también fuera de constructores (guards, interceptores, funciones).

\paragraph{Scopes útiles}
\texttt{providedIn: 'root'} = singleton global. \texttt{providedIn: 'any'} = una instancia por inyector (útil en \emph{lazy}). \texttt{providers: [...]} en un componente = instancias por vista (evita fugas entre rutas).

\begin{lstlisting}[style=tsstyle, caption={Tokens, scopes y inject()}, label={lst:di}]
// tokens.ts
import { InjectionToken, inject, DestroyRef } from '@angular/core';

export type ApiConfig = Readonly<{ baseUrl: string; timeoutMs: number }>; 
export const API_CONFIG = new InjectionToken<ApiConfig>('API_CONFIG');

export function useAbortable<T>(subscribe: (signal: AbortSignal) => Promise<T>) {
  const destroyRef = inject(DestroyRef);
  const ac = new AbortController();
  destroyRef.onDestroy(() => ac.abort());
  return subscribe(ac.signal);
}
\end{lstlisting}

\begin{warnbox}
No acoples componentes a servicios concretos: usa \texttt{InjectionToken} para configuraciones y facilita los tests.
\end{warnbox}

\section{Componentes profesionales}
\paragraph{Cómo pensar un componente}
Un componente debería: (1) recibir datos por \texttt{input()} y emitir acciones por \texttt{output()}, (2) gestionar sólo estado de UI local con Signals, (3) delegar I/O a servicios.

\subsection{Inputs/Outputs, Signals y Change Detection}
\paragraph{Qué hace cada pieza}
- \textbf{Inputs/Outputs}: API pública del componente. \texttt{input<number>(1)} tipa y da valor por defecto. \texttt{output<number}()> emite eventos tipados.
- \textbf{Signals}: \texttt{signal} = estado, \texttt{computed} = derivado, \texttt{effect} = efectos secundarios (persistencia, analytics).
- \textbf{OnPush}: evita renders innecesarios; cambia la vista cuando cambian inputs, signals o hay eventos.

\begin{lstlisting}[style=tsstyle, caption={Componente con Signals, OnPush y efecto de sincronización}]
import { Component, signal, computed, effect, input, output } from '@angular/core';

@Component({
  selector: 'app-counter',
  standalone: true,
  templateUrl: './counter.html',
  changeDetection: 0 // ChangeDetectionStrategy.OnPush
})
export class CounterComponent {
  step = input<number>(1);     // nuevo input con tipado fuerte
  changed = output<number>();  // nuevo output funcional

  count = signal(0);
  doubled = computed(() => this.count() * 2);

  constructor() {
    effect(() => {
      const c = this.count();
      // sincroniza con almacenamiento, analytics, etc.
      localStorage.setItem('count', String(c));
    });
  }

  inc() {
    this.count.update(c => c + this.step());
    this.changed.emit(this.count());
  }
}
\end{lstlisting}

\begin{lstlisting}[style=htmlstyle, caption={Plantilla moderna con @if/@for y @defer}]
<div class="panel">
  <h2>Contador: {{ count() }} (x2: {{ doubled() }})</h2>

  @if (count() === 0) {
    <p>Empieza a contar</p>
  } @else {
    <p>Llevas {{ count() }} clicks</p>
  }

  <button (click)="inc()">+{{ step() }}</button>

  @defer (on idle; prefetch on viewport) {
    <lazy-stats [value]="count()"/>
  } @placeholder {
    <p>Cargando estadísticas…</p>
  }
</div>
\end{lstlisting}

\begin{warnbox}
No mutar arrays/objetos en sitio. Usa \texttt{update} con copias (\texttt{[...prev, nuevo]}) para mantener OnPush eficiente.
\end{warnbox}

\subsection{Proyección, referencias y hooks}
\paragraph{Cuándo usar content projection}
Cuando un componente \emph{contenedor} (card, modal) debe renderizar contenido arbitrario del padre. Usa \texttt{ng-content} o \texttt{@ContentChild} para plantillas.

\begin{lstlisting}[style=tsstyle, caption={Content projection y limpieza reactiva}]
import { Component, ContentChild, TemplateRef } from '@angular/core';
import { takeUntilDestroyed } from '@angular/core/rxjs-interop';
import { interval } from 'rxjs';

@Component({ selector: 'app-card', standalone: true, templateUrl: './card.html' })
export class CardComponent {
  @ContentChild(TemplateRef) bodyTpl?: TemplateRef<unknown>;

  ngOnInit() {
    interval(1000)
      .pipe(takeUntilDestroyed())
      .subscribe(() => console.log('tick'));
  }
}
\end{lstlisting}

\section{Formularios tipados y validaciones}
\paragraph{Elegir estrategia}
- \textbf{Reactive Forms}: más predecibles, testables y tipables. Preferidos en apps medianas/grandes.
- \textbf{Template-driven}: válidos para formularios muy simples.

\paragraph{Mostrar errores de forma accesible}
Lee \texttt{touched}/\texttt{dirty} para decidir cuándo mostrar errores. Añade \texttt{aria-invalid} y \texttt{aria-describedby}.

\subsection{Reactive Forms}
\begin{lstlisting}[style=tsstyle, caption={FormGroup tipado y validadores personalizados}]
import { Component } from '@angular/core';
import { FormControl, FormGroup, Validators, NonNullableFormBuilder } from '@angular/forms';

interface ProfileForm {
  name: FormControl<string>;
  email: FormControl<string>;
}

@Component({ selector: 'app-profile', standalone: true, templateUrl: './profile.html' })
export class ProfileComponent {
  readonly form = new FormGroup<ProfileForm>({
    name: new FormControl('', { nonNullable: true, validators: [Validators.required] }),
    email: new FormControl('', { nonNullable: true, validators: [Validators.email] })
  });
}
\end{lstlisting}

\begin{lstlisting}[style=tsstyle, caption={Validador asíncrono con cache y cancelación}]
import { AbstractControl, AsyncValidatorFn, ValidationErrors } from '@angular/forms';
import { timer, map, switchMap, of } from 'rxjs';

export function emailTaken(check: (email: string) => Promise<boolean>): AsyncValidatorFn {
  return (ctrl: AbstractControl) => {
    const email = ctrl.value as string;
    if (!email) return of(null);
    return timer(300).pipe( // debounce
      switchMap(() => check(email)),
      map(isTaken => (isTaken ? { emailTaken: true } as ValidationErrors : null))
    );
  };
}
\end{lstlisting}

\section{Routing funcional, lazy y resolución de datos}
\paragraph{Ideas clave}
Rutas como datos (\texttt{Routes}); \texttt{loadComponent}/\texttt{loadChildren} para \emph{lazy}; Guards controlan acceso; Resolvers traen datos antes de entrar; usa \texttt{data} para metadatos y \texttt{ActivatedRoute} para leer params.

\begin{lstlisting}[style=tsstyle, caption={Definición de rutas standalone con guards funcionales}]
// app.routes.ts
import { Routes, provideRouter, CanActivateFn, ResolveFn } from '@angular/router';
import { inject } from '@angular/core';
import { AuthService } from './auth.service';

const canActivateAuth: CanActivateFn = () => inject(AuthService).isLoggedIn();
const resolveUser: ResolveFn<User> = () => inject(AuthService).loadCurrentUser();

export const routes: Routes = [
  { path: '', loadComponent: () => import('./home.component').then(m => m.HomeComponent) },
  { path: 'admin', canActivate: [canActivateAuth],
    loadChildren: () => import('./admin/routes').then(m => m.ADMIN_ROUTES),
    resolve: { user: resolveUser }
  },
  { path: '**', redirectTo: '' }
];
\end{lstlisting}

\begin{tipbox}
Evita lógica de negocio en guards/resolvers. Delega en servicios y mantén funciones puras y testeables.
\end{tipbox}

\section{HTTP profesional: interceptores, errores, caché}
\paragraph{Patrones esenciales}
- Interceptor de \textbf{auth}: añade token y reacciona a 401 (logout/refresh).
- \textbf{Retry con backoff}: reintenta fallos transitorios sin saturar el servidor.
- \textbf{Caché}: para GET idempotentes en listas/constantes.

\begin{lstlisting}[style=tsstyle, caption={Interceptor de auth y reintento exponencial}]
import { HttpInterceptorFn, HttpRequest, HttpHandlerFn } from '@angular/common/http';
import { inject } from '@angular/core';
import { catchError, retry, timer, mergeMap, throwError } from 'rxjs';

export const authInterceptor: HttpInterceptorFn = (req: HttpRequest<unknown>, next: HttpHandlerFn) => {
  const token = inject(TokenService).get();
  const cloned = token ? req.clone({ setHeaders: { Authorization: `Bearer ${token}` }}) : req;
  return next(cloned).pipe(
    retry({ count: 2, delay: (e, i) => timer(2 ** i * 200) }),
    catchError(err => {
      if (err.status === 401) inject(AuthService).logout();
      return throwError(() => err);
    })
  );
};
\end{lstlisting}

\begin{lstlisting}[style=tsstyle, caption={Interceptor de caché simple por URL y método}]
const cache = new Map<string, unknown>();
export const cacheInterceptor: HttpInterceptorFn = (req, next) => {
  if (req.method !== 'GET') return next(req);
  const key = req.urlWithParams;
  if (cache.has(key)) return of(new HttpResponse({ status: 200, body: cache.get(key) })).pipe(delay(0));
  return next(req).pipe(tap((event) => {
    if (event instanceof HttpResponse) cache.set(key, event.body);
  }));
};
\end{lstlisting}

\begin{warnbox}
No caches peticiones mutantes (POST/PUT/DELETE) y recuerda invalidar caché tras cambios.
\end{warnbox}

\section{Estado: RxJS, Signals y NgRx}
\paragraph{Cuándo usar cada uno}
- \textbf{Signals}: estado local de un componente o pequeñas vistas.
- \textbf{RxJS}: flujos asíncronos (eventos, websockets, composición de peticiones) y \emph{streams} compartidos.
- \textbf{NgRx}: estado global complejo, múltiples actores, time-travel/debugging, reglas de negocio compartidas.

\subsection{RxJS en componentes}
\paragraph{Elegir operador correctamente}
\texttt{switchMap} cancela lo anterior (búsqueda), \texttt{concatMap} ordena (cola), \texttt{mergeMap} paraleliza (orden no importa), \texttt{exhaustMap} evita dobles envíos.

\begin{lstlisting}[style=tsstyle, caption={Elección correcta entre mergeMap/switchMap/concatMap/exhaustMap}]
// Regla mental:
// - switchMap: cancelar lo anterior (búsqueda incremental)
// - concatMap: en cola (operaciones ordenadas)
// - mergeMap: en paralelo (cuando el orden no importa)
// - exhaustMap: ignora hasta terminar (evita doble submit)
\end{lstlisting}

\subsection{Signals para estado local}
\paragraph{Regla práctica}
Deriva todo lo calculable con \texttt{computed} y mantén \texttt{signal} como fuente de verdad mínima.

\begin{lstlisting}[style=tsstyle, caption={Composición de estado con signal/computed/effect}]
export class TodosStore {
  readonly todos = signal<ReadonlyArray<Todo>>([]);
  readonly query = signal('');
  readonly filtered = computed(() => this.todos().filter(t => t.title.includes(this.query())));
  add(title: string) { this.todos.update(arr => [...arr, { id: crypto.randomUUID(), title, done: false }]); }
}
\end{lstlisting}

\subsection{NgRx moderno}
\paragraph{Estructura mental}
Acciones describen \emph{qué pasó}, reducers calculan \emph{nuevo estado}, selectors exponen \emph{lecturas eficientes}. Efectos coordinan I/O.

\begin{lstlisting}[style=tsstyle, caption={Slice con createFeature y efectos tipados}]
import { createActionGroup, props, createFeature, createReducer, on, createSelector } from '@ngrx/store';

export const UsersActions = createActionGroup({ source: 'Users', events: {
  'Load': props<void>(),
  'Loaded Success': props<{ users: User[] }>(),
  'Loaded Failure': props<{ error: unknown }>(),
}});

export type UsersState = Readonly<{ users: User[]; loading: boolean; error?: unknown }>; 
const initialState: UsersState = { users: [], loading: false };

export const usersFeature = createFeature({
  name: 'users',
  reducer: createReducer(initialState,
    on(UsersActions.load, (s) => ({ ...s, loading: true })),
    on(UsersActions.loadedSuccess, (s, { users }) => ({ ...s, loading: false, users })),
    on(UsersActions.loadedFailure, (s, { error }) => ({ ...s, loading: false, error })),
  )
});

export const { name, reducer, selectUsers, selectLoading } = usersFeature;
\end{lstlisting}

\section{Rendimiento y escalabilidad}
\paragraph{Qué optimizar primero}
- Cambia a \textbf{OnPush}, usa \texttt{trackBy} y evita trabajo pesado en plantillas.
- Divide rutas y usa \texttt{@defer} para cargar tarde lo no crítico.
- Considera \emph{zoneless} cuando controles el flujo con Signals y eventos propios.

\begin{itemize}[leftmargin=*]
  \item \textbf{OnPush} + inmutabilidad + \texttt{trackBy} en \texttt{@for}.
  \item Evita pipes pesados en plantilla; prefiera \texttt{computed} o pre-cálculo.
  \item Divide rutas y usa \texttt{@defer} para no bloquear TTI.
  \item Considera \textbf{zoneless}: \texttt{bootstrapApplication(App, { providers: [provideZoneChangeDetection({ eventCoalescing: true, runCoalescing: true })] })} o \texttt{ngZone: 'noop'} y dispara CD manual con Signals.
  \item \textbf{SSR + Hydration}: mejora SEO y LCP.
\end{itemize}

\begin{lstlisting}[style=tsstyle, caption={Lista grande con @for y trackBy}]
// template
@for (item of items(); track item.id) {
  <app-row [item]="item"/>
}
\end{lstlisting}

\section{SSR, Hydration y despliegue}
\paragraph{Cómo encajan}
SSR renderiza HTML en el servidor para entregar contenido rápido y rastreable. En cliente, \emph{hydration} conecta eventos sin re-render completo. Protege el uso de APIs del navegador con \texttt{isPlatformBrowser}.

\begin{itemize}[leftmargin=*]
  \item Genera SSR: \texttt{ng add @nguniversal/express-engine}
  \item Verifica rutas \emph{lazy} y uso de APIs del browser en \texttt{isPlatformBrowser}.
  \item Habilita hidración con \texttt{provideClientHydration()} en el bootstrap.
\end{itemize}

\begin{lstlisting}[style=tsstyle, caption={Bootstrap con Hydration}]
import { provideClientHydration } from '@angular/platform-browser';
bootstrapApplication(AppComponent, { providers: [ provideClientHydration() ] });
\end{lstlisting}

\section{Testing eficaz}
\paragraph{Estrategia mínima}
- Unit: componentes standalone con TestBed, servicios con dobles de prueba.
- Integración: prueba flujo entre componentes/servicios.
- E2E: escenarios críticos (login, compra). Usa \emph{data-testid} y retrys.

\begin{lstlisting}[style=tsstyle, caption={Prueba de componente standalone}]
import { TestBed } from '@angular/core/testing';
import { AppComponent } from './app.component';

describe('AppComponent', () => {
  it('renderiza título', async () => {
    await TestBed.configureTestingModule({ imports: [AppComponent] }).compileComponents();
    const fixture = TestBed.createComponent(AppComponent);
    fixture.detectChanges();
    expect(fixture.nativeElement.querySelector('h1').textContent).toContain('Hola');
  });
});
\end{lstlisting}

\section{Accesibilidad, i18n y seguridad}
\paragraph{Buenas prácticas inmediatas}
Añade roles/ARIA, gestiona el foco tras diálogos y navegación, evita \texttt{[innerHTML]} sin sanitizar, marca cadenas con \texttt{$localize} para traducción.

\begin{itemize}[leftmargin=*]
  \item Usa roles/ARIA, focus management y \texttt{cdkTrapFocus}.
  \item i18n con \texttt{@angular/localize}; evita concatenaciones dinámicas.
  \item Sanear HTML: \texttt{DomSanitizer}; evita \texttt{bypassSecurityTrustXxx} salvo casos justificados.
\end{itemize}

\section{Herramientas y calidad}
\paragraph{Cómo empezar}
Activa ESLint con \texttt{angular-eslint}, configura Prettier y añade un pipeline de CI con build y tests. Analiza bundles antes de desplegar.

\begin{itemize}[leftmargin=*]
  \item ESLint con reglas \emph{angular-eslint}; \texttt{no-implicit-any}, \texttt{no-floating-promises}.
  \item Prettier y formateo consistente.
  \item Analiza bundles: \texttt{ng build --configuration production --stats-json} + \texttt{source-map-explorer}.
  \item Monorepo: considera Nx para escalabilidad.
\end{itemize}

\section{Patrones y anti-patrones}
\begin{refbox}
Patrones: Presentational/Container, Facade para estado, Smart HTTP Service, Component Store para vistas complejas, Signals para estado local.
\end{refbox}

\begin{warnbox}
Anti-patrones: servicios \emph{god}, mutaciones silenciosas, suscribir en servicios sin exponer Observables, abusar de \texttt{any}, lógica de negocio en plantillas.
\end{warnbox}

\section{Checklist de revisión senior}
\begin{enumerate}[leftmargin=*]
  \item Tipado estricto activado; sin \texttt{any}/\texttt{@ts-ignore}.
  \item OnPush y trackBy en listas.
  \item División de rutas y \texttt{@defer} aplicado.
  \item Interceptores para auth/errores; reintentos y \emph{backoff}.
  \item Limpieza con \texttt{takeUntilDestroyed}; cero fugas.
  \item Accesibilidad validada y \texttt{aria-*} correcto.
  \item Tests unitarios y de integración para piezas críticas.
\end{enumerate}

% ------------------------------------------------------
% Ruta de aprendizaje y refuerzo de fundamentos
% ------------------------------------------------------
\section{Ruta de aprendizaje a experto (temario por niveles)}
Esta ruta está pensada para construir bases sólidas, dominar el framework y entender sus internals. Muchas secciones enlazan con contenidos previos del documento.

\subsection{1. Fundamentos del ecosistema Angular}
\begin{itemize}[leftmargin=*]
  \item SPA: una única carga de documento, navegación cliente, estado en memoria, API vía HTTP/WS.
  \item Angular vs React vs Vue: Angular es un \emph{framework} completo (router, DI, formularios, i18n); React y Vue son \emph{libraries} que se amplían con ecosistema.
  \item TypeScript base: tipos, genéricos, union/intersection, narrowing, decoradores, módulos ES.
  \item Node/npm: scripts, semver, workspaces, \texttt{npx}.
  \item Modularidad y arquitectura: separación por \textbf{feature}, \textbf{core}, \textbf{shared}; patrones Presentational/Container.
\end{itemize}

\subsection{2. Estructura básica de un proyecto}
\begin{itemize}[leftmargin=*]
  \item CLI: crear, generar artefactos, \texttt{ng generate}, \texttt{ng build}, \texttt{ng test}.
  \item Carpetas: \texttt{src/app}, \texttt{assets}, \texttt{environments}. En apps modernas, \textbf{standalone} reemplaza \texttt{NgModule}.
  \item Componentes: \texttt{@Component}, plantilla, estilos, ciclo de vida.
\end{itemize}

\subsection{3. Templates y data binding}
\begin{itemize}[leftmargin=*]
  \item Interpolación, property/event binding, two-way con \texttt{[(ngModel)]} (formularios de plantilla).
  \item Directivas estructurales y control de flujo moderno: \texttt{@if}, \texttt{@for}, \texttt{@switch}.
  \item Directivas de atributo: \texttt{[class]}, \texttt{[style]}, \texttt{[ngClass]}, \texttt{[ngStyle]}.
\end{itemize}

\subsection{4. Comunicación entre componentes}
\begin{itemize}[leftmargin=*]
  \item \texttt{@Input()} y \texttt{@Output()} o APIs \texttt{input()/output()}.
  \item \texttt{ViewChild}/\texttt{ContentChild}/\texttt{ViewChildren}.
  \item Servicios compartidos como canal de eventos/estado.
\end{itemize}

\subsection{5. Servicios e inyección de dependencias}
\begin{itemize}[leftmargin=*]
  \item \texttt{@Injectable}, \texttt{providedIn}, \textbf{scope} raíz vs componente.
  \item Jerarquía de inyectores, \texttt{inject()} y \texttt{InjectionToken} (ver Listado~\ref{lst:di}).
\end{itemize}

\subsection{6. Ruteo y navegación}
\begin{itemize}[leftmargin=*]
  \item Definición de \texttt{Routes}, navegación, parámetros, rutas hijas.
  \item Lazy loading con \texttt{loadComponent}/\texttt{loadChildren}. Guards y resolvers.
\end{itemize}

\subsection{7. Formularios}
\begin{itemize}[leftmargin=*]
  \item Template-driven y Reactive Forms; validaciones sincronas/asincronas; estados touched/dirty/pristine.
  \item Mensajería de errores amigable y accesible.
\end{itemize}

\subsection{8. Observables y RxJS}
\begin{itemize}[leftmargin=*]
  \item Observables, suscripción, cancelación; operadores clave (\texttt{map}, \texttt{filter}, \texttt{switchMap}, \texttt{catchError}).
  \item \texttt{Subject}/\texttt{BehaviorSubject}/\texttt{ReplaySubject}; composición y \texttt{shareReplay}.
\end{itemize}

\subsection{9. HTTP Client}
\begin{itemize}[leftmargin=*]
  \item \texttt{HttpClient}: \texttt{get/post/put/delete}, cabeceras, parámetros, \texttt{observe} vs \texttt{responseType}.
  \item Interceptores, manejo de errores, JWT y refresh tokens.
\end{itemize}

\subsection{10. Arquitectura avanzada}
\begin{itemize}[leftmargin=*]
  \item División por dominio: feature folders, \texttt{Core} y \texttt{Shared}. Facades y servicios HTTP delgados.
  \item Smart vs Dumb Components; \textbf{Signals} para estado local, NgRx/Component Store para vistas complejas.
\end{itemize}

\subsection{11. Gestión de estado}
\begin{itemize}[leftmargin=*]
  \item NgRx: \emph{actions, reducers, selectors, effects}. Alternativa \textbf{Signals} (Angular 16+), ver secciones previas.
\end{itemize}

\subsection{12. Rendimiento y optimización}
\begin{itemize}[leftmargin=*]
  \item \texttt{OnPush}, inmutabilidad, \texttt{trackBy}, \texttt{@defer}, zoneless, SSR + hydration.
\end{itemize}

\subsection{13. Testing}
\begin{itemize}[leftmargin=*]
  \item Unit con TestBed; mocks y espías; pruebas de servicios/pipes/directivas.
  \item E2E: \textbf{Cypress/Playwright} (Protractor está deprecado).
\end{itemize}

\subsection{14. Despliegue y DevOps}
\begin{itemize}[leftmargin=*]
  \item \texttt{ng build --configuration production}, entornos, variables, \texttt{docker} y CDNs.
  \item CI/CD con GitHub Actions (build, test, deploy).
\end{itemize}

\subsection{15. Ecosistema avanzado}
\begin{itemize}[leftmargin=*]
  \item Angular Material y CDK, Animaciones, Directivas/Pipes personalizadas, Angular Elements, WebSockets.
\end{itemize}

\subsection{16. Buenas prácticas}
\begin{itemize}[leftmargin=*]
  \item Convenciones, reutilización, arquitectura limpia, documentación, \texttt{eslint} y \texttt{prettier}.
\end{itemize}

% ------------------------------------------------------
% TypeScript avanzado
% ------------------------------------------------------
\section{TypeScript avanzado y tipado estricto en Angular}
\subsection{Tipos condicionales, utilitarios y ``type-level programming''}
\begin{lstlisting}[style=tsstyle, caption={Condicionales, infer y mapeados útiles}]
// Type guard para APIs estrictas
export function isUser(x: unknown): x is User {
  return !!x && typeof x === 'object' && 'id' in (x as any) && 'email' in (x as any);
}

// Condicional con infer: extraer payload de una Action NgRx
export type Action = { type: string; payload?: unknown };
export type PayloadOf<A extends Action> = A extends { payload: infer P } ? P : never;

// Mapeados para construir formularios tipados desde un modelo
export type Formify<T> = { [K in keyof T]-?: FormControl<T[K]> };

type UserModel = { name: string; age: number };
export type UserForm = FormGroup<Formify<UserModel>>;
\end{lstlisting}

\subsection{Decoradores y metadatos}
\begin{itemize}[leftmargin=*]
  \item En Angular moderno, los decoradores se usan a nivel \emph{estático}; el compilador genera metadatos internos de Ivy en tiempo de compilación (AOT/Ivy), por lo que no se requiere \texttt{reflect-metadata} en producción.
  \item Evita usar decoradores de clase personalizados en runtime; prefiere funciones puras y \texttt{InjectionToken} para configuración.
\end{itemize}

\subsection{Chequeo de tipos en plantillas}
\begin{itemize}[leftmargin=*]
  \item El \emph{Template Type Checker} valida expresiones en HTML (strictTemplates). Actívalo siempre.
  \item Usa \texttt{as const}, narrowings y \texttt{readonly} para ayudar al compilador.
\end{itemize}

\begin{lstlisting}[style=tsstyle, caption={Tokens tipados y configuración segura}]
export interface FeatureFlags { readonly enableBeta: boolean; }
export const FEATURE_FLAGS = new InjectionToken<FeatureFlags>('FEATURE_FLAGS');

bootstrapApplication(AppComponent, { providers: [
  { provide: FEATURE_FLAGS, useValue: { enableBeta: true } satisfies FeatureFlags }
]});
\end{lstlisting}

% ------------------------------------------------------
% Internals: cómo funciona Angular por debajo
% ------------------------------------------------------
\section{Angular por debajo del capó (internals esenciales)}
\subsection{Change Detection: de Zone.js a Signals}
\begin{itemize}[leftmargin=*]
  \item \textbf{Zone.js}: parchea eventos/microtareas; tras eventos HTTP/DOM, marca vistas para comprobar. \texttt{OnPush} limita la propagación a inputs cambiados y eventos.
  \item \textbf{Zoneless + Signals}: los \emph{signals} mantienen un grafo de dependencias; \texttt{computed} recalcula cuando cambian sus fuentes y las vistas leen señales con \emph{auto-dirty-check} eficiente.
  \item APIs útiles: \texttt{markDirty}, \texttt{ChangeDetectorRef.detectChanges()}, \texttt{runOutsideAngular} para evitar trabajo extra.
\end{itemize}

\subsection{Inyección de dependencias en Ivy}
\begin{itemize}[leftmargin=*]
  \item Cada componente/directiva se asocia a vistas internas (\emph{LView/TView}). La resolución de \texttt{inject()} sube por el \emph{árbol de inyectores} hasta encontrar un proveedor o delega al inyector raíz.
  \item \texttt{providedIn: 'any'} crea un singleton por inyector; \texttt{providers: [...]} en componentes crea instancias por vista.
\end{itemize}

\subsection{Compilación de plantillas (Ivy Instructions)}
\begin{itemize}[leftmargin=*]
  \item La plantilla se transforma a instrucciones como instrucciones de creación/actualización separadas, optimizando el trabajo en cada ciclo de CD.
  \item El control de flujo \texttt{@if/@for} compila a bloques eficientes con \emph{block synthesis} y \emph{dom diffing} mínimo.
\end{itemize}

\begin{lstlisting}[style=tsstyle, caption={Esquema simplificado de plantilla compilada}]
// Pseudocódigo: create/update
function Template(rf: RenderFlags, ctx: Ctx) {
  if (rf & 1) { /* create phase */
    // elementStart('h1');
    // text();
    // elementEnd();
  }
  if (rf & 2) { /* update phase */
    // textInterpolate(ctx.title);
  }
}
\end{lstlisting}

\subsection{Hydration y SSR}
\begin{itemize}[leftmargin=*]
  \item SSR genera HTML inicial; en cliente, Angular \emph{hidrata}: empareja nodos y adjunta listeners sin re-render completo.
  \item Evita usar APIs de navegador durante SSR; protege con \texttt{isPlatformBrowser}.
\end{itemize}

\subsection{Builder y empaquetado}
\begin{itemize}[leftmargin=*]
  \item El builder actual usa \textbf{esbuild} para velocidad; \emph{treeshaking}, \emph{code-splitting}, CSS inlining y optimización de imágenes.
  \item Habilita \texttt{optimization}, \texttt{budgets} y \texttt{stats-json} para auditar tamaño.
\end{itemize}

% ------------------------------------------------------
% Plan de estudio sugerido
% ------------------------------------------------------
\section{Plan de estudio recomendado (6 semanas)}
\begin{enumerate}[leftmargin=*]
  \item Semana 1: Fundamentos TS, CLI, estructura y componentes. Labs: 3 componentes + comunicación básica.
  \item Semana 2: Templates, control de flujo moderno, formularios template-driven y Reactive básicos.
  \item Semana 3: HTTP, interceptores, RxJS esencial (switchMap, error handling), arquitectura por features.
  \item Semana 4: Estado con Signals y NgRx intro; performance (OnPush, trackBy); routing avanzado y lazy.
  \item Semana 5: SSR + hydration, zoneless, accesibilidad e i18n, animaciones y Material/CDK.
  \item Semana 6: Testing unitario y E2E, seguridad (JWT/CSRF), CI/CD y despliegue; auditoría de bundles.
\end{enumerate}

\begin{tipbox}
Estudia construyendo. Define un mini-producto (lista de tareas con auth + SSR) e incrementa complejidad cada semana aplicando los conceptos.
\end{tipbox}

% ------------------------------------------------------
% Ecosistema avanzado Angular: Material, CDK, Animaciones, PWA, i18n, WS
% ------------------------------------------------------
\section{Material, CDK y diseño de UI}
\subsection{Angular Material (standalone)}
\begin{lstlisting}[style=tsstyle, caption={Uso de Material en standalone}]
import { Component } from '@angular/core';
import { MatButtonModule } from '@angular/material/button';
import { MatIconModule } from '@angular/material/icon';

@Component({
  selector: 'app-toolbar',
  standalone: true,
  imports: [MatButtonModule, MatIconModule],
  template: `
    <button mat-raised-button color="primary"><mat-icon>save</mat-icon> Guardar</button>
  `
})
export class ToolbarComponent {}
\end{lstlisting}

\subsection{CDK Overlay y Portal}
\begin{lstlisting}[style=tsstyle, caption={Overlay básico con CDK}]
import { Overlay, OverlayRef } from '@angular/cdk/overlay';
import { ComponentPortal } from '@angular/cdk/portal';

const overlayRef: OverlayRef = inject(Overlay).create();
overlayRef.attach(new ComponentPortal(MyPopupComponent));
\end{lstlisting}

\section{Animaciones}
\begin{lstlisting}[style=tsstyle, caption={Trigger con estados y transiciones}]
import { trigger, state, style, transition, animate } from '@angular/animations';

@Component({
  // ...
  animations: [
    trigger('fade', [
      state('in', style({ opacity: 1 })),
      transition(':enter', [ style({ opacity: 0 }), animate('200ms ease-out') ]),
      transition(':leave', [ animate('150ms ease-in', style({ opacity: 0 })) ])
    ])
  ]
})
export class Panel {}
\end{lstlisting}

\section{PWA y Service Worker}
\begin{itemize}[leftmargin=*]
  \item Añade PWA: \texttt{ng add @angular/pwa}. Configura \texttt{ngsw-config.json} y \texttt{provideServiceWorker()}.
\end{itemize}
\begin{lstlisting}[style=tsstyle, caption={Bootstrap con Service Worker}]
import { provideServiceWorker } from '@angular/service-worker';
bootstrapApplication(AppComponent, {
  providers: [ provideServiceWorker('ngsw-worker.js', { enabled: true }) ]
});
\end{lstlisting}

\section{i18n y localización}
\begin{lstlisting}[style=tsstyle, caption={Marcado con \$localize y runtime i18n}]
const title = $localize`Hola, ${name}:bienvenido!`;
\end{lstlisting}
\begin{itemize}[leftmargin=*]
  \item Extrae mensajes: \texttt{ng extract-i18n} y compila con \texttt{@angular/localize}.
\end{itemize}

\section{WebSockets y tiempo real}
\begin{lstlisting}[style=tsstyle, caption={Servicio con WebSocketSubject}]
import { webSocket } from 'rxjs/webSocket';
@Injectable({ providedIn: 'root' })
export class WsService {
  private socket$ = webSocket<string>('wss://example.com/ws');
  messages$ = this.socket$.asObservable();
  send(msg: string) { this.socket$.next(msg); }
}
\end{lstlisting}

\section{Web Workers}
\begin{lstlisting}[style=tsstyle, caption={Creación y comunicación con Worker}]
const worker = new Worker(new URL('./heavy.worker', import.meta.url));
worker.postMessage({ task: 'calc' });
worker.onmessage = (e) => console.log(e.data);
\end{lstlisting}

\section{Seguridad en el navegador}
\begin{itemize}[leftmargin=*]
  \item XSS: el binding Angular es seguro por defecto; \texttt{[innerHTML]} requiere sanitización. Evita \texttt{bypassSecurityTrustXxx} salvo casos controlados.
  \item CSP: usa nonces/hashes y evita \texttt{unsafe-inline}. Verifica 3rd-party scripts.
\end{itemize}

\section{Estrategias de precarga y performance de rutas}
\begin{lstlisting}[style=tsstyle, caption={PreloadAllModules y precarga personalizada}]
import { withPreloading, PreloadAllModules } from '@angular/router';
provideRouter(routes, withPreloading(PreloadAllModules));
\end{lstlisting}

\begin{lstlisting}[style=tsstyle, caption={Precarga condicional por data}]
export const routes: Routes = [
  { path: 'admin', loadChildren: () => import('./admin/routes').then(m => m.ADMIN_ROUTES), data: { preload: true } }
];

@Injectable({ providedIn: 'root' })
export class DataPreloadingStrategy implements PreloadingStrategy {
  preload(route: Route, load: () => Observable<unknown>) {
    return route.data?.['preload'] ? load() : of(null);
  }
}
\end{lstlisting}

% ------------------------------------------------------
% Monorepo y DX
% ------------------------------------------------------
\section{Monorepo y DX}
\begin{itemize}[leftmargin=*]
  \item Nx: \emph{generators}, \emph{executors}, caché distribuida, grafo de dependencias.
  \item Auditoría: \texttt{ng build --stats-json} + \texttt{source-map-explorer}.
\end{itemize}

\end{document}
